%=====================================================================
% SECTION 5: RESULTS AND DISCUSSION  (skeleton; not part of paper.tex)
% Structure and paragraph plan: reports/section5_structure.md
% Numbers: reports/section5_facts.md and reports/section5/*.csv
%   (scripts/section5/s5_facts.py). Every % comment below is a paragraph
%   to write; numbers in comments are copied from the facts sheet.
% Figure numbers assume paper Figures 1-3 = area, gauges, model.
% If Figure 6 moves to the SI, Figures 7-10 become 6-9.
%=====================================================================
\section{Results and Discussion}
\label{sec:results}

% P1 Comparisons and status rule. M0-LSTM (RQ1); M0-M2, M1-M4 (RQ2);
%    M0-M1 (RQ3); NWM reference only. n = 101 gauged, 98 random hold-out,
%    29 blocked; no gauged M4. Status: equivalent (Holm TOST p<0.05 at the
%    margin) / different (Holm Wilcoxon p<0.05, interval excludes 0) /
%    undetermined. Mean paired difference in days 1-5 NSE (median beside
%    it), 90% gauge-bootstrap intervals; one training run per model, so
%    intervals are across gauges. All contrasts: Table~\ref{tab:contrasts}.

%---------------------------------------------------------------------
\subsection{Skill Across Lead Times and Protocols}
\label{sec:res-skill}
% RESULTS
% P1 Answer + Table 3: days 1-5 median NSE M0 0.430/0.385/0.305,
%    LSTM 0.428/0.367/0.280, NWM 0.255/0.255/0.183 (gauged/random/blocked).
%    Skill falls with hold-out strictness for every model. Figure 4 top row:
%    all 40 leads; curves converge with lead.
% P2 KGE vs NSE: blocked KGE LSTM 0.381, NWM 0.385, M0 0.334; alpha M0 0.530
%    vs LSTM 0.626 (blocked); gauged KGE M0 0.526 vs LSTM 0.481.
%    NWM better than M0 at 6/101, 17/98, 5/29 gauges.
% P3 Peaks: top-5% bias -0.37 to -0.48 for every model and protocol; peaks
%    within 24 h timed at 3 h (M0) vs 6 h (LSTM, NWM), blocked 1.5 vs 6 h.
%    Example event: Figure S3.
% DISCUSSION
% P4 NWM is nudged (reference, not competitor). Convergence with lead:
%    precipitation-forecast error dominates \cite{Cloke2009,Wood2008}.
% P5 NSE rewards damped forecasts \cite{Gupta2009}: lower alpha of M0,
%    peak underestimation, relevance for flood warning.

\begin{figure}
\figfile{figures/fig4_skill_lead.pdf}{\textwidth}
\caption{Skill against lead time, 6 to 240~h, for the reach-scale model (M0), the LSTM, and the NWM forecast in the (a) gauged ($n = 101$), (b) random hold-out ($n = 98$), and (c) blocked ($n = 29$) protocols. Top: median NSE over gauges. Bottom: paired difference M0 $-$ LSTM, mean over gauges with 90\% gauge-bootstrap band. Days 1--5 are shaded.}
\label{fig:skill}
\end{figure}

\begin{table}
\caption{Median NSE and KGE, Days 1--5, by Protocol and Model}
\label{tab:skill}
\footnotesize
\setlength{\parindent}{0pt}
\renewcommand{\arraystretch}{1.2}
\begin{tabular}{@{}llrrr@{}}
\toprule
Protocol & Model & $n$ & NSE & KGE \\
\midrule
Gauged & GNN & 101 & \textbf{0.430} & \textbf{0.526} \\
 & GNN-noIns & 101 & 0.428 & \textbf{0.526} \\
 & GNN-noMP & 101 & 0.426 & 0.517 \\
 & GNN-Instat & --- & \multicolumn{2}{c}{not trained} \\
 & LSTM & 101 & 0.428 & 0.481 \\
 & NWM & 101 & 0.255 & 0.443 \\
\addlinespace
Random hold-out & GNN & 98 & 0.385 & 0.429 \\
 & GNN-noIns & 98 & 0.383 & 0.428 \\
 & GNN-noMP & 98 & \textbf{0.386} & 0.421 \\
 & GNN-Instat & 98 & 0.379 & 0.421 \\
 & LSTM & 98 & 0.367 & 0.396 \\
 & NWM & 98 & 0.255 & \textbf{0.448} \\
\addlinespace
Blocked basins & GNN & 29 & 0.305 & 0.334 \\
 & GNN-noIns & 29 & 0.297 & 0.334 \\
 & GNN-noMP & 29 & \textbf{0.322} & 0.341 \\
 & GNN-Instat & 29 & 0.307 & 0.324 \\
 & LSTM & 29 & 0.280 & 0.381 \\
 & NWM & 29 & 0.183 & \textbf{0.385} \\
\bottomrule
\end{tabular}
\begin{flushleft}
\footnotesize \emph{Note.} Median over the $n$ gauges scored by every model of the protocol, per lead, averaged over the leads 24, 48, 72, 96 and 120~h. Bold: the highest value per protocol. Differences below 0.02 NSE are not distinguishable (equivalence margin, Table~S6). GNN-Instat was not trained for the gauged protocol. NWM: operational forecast, nudged toward observations at gauges, a reference only.
\end{flushleft}
\end{table}


%---------------------------------------------------------------------
\subsection{Reach-Scale Model Versus the Regional LSTM (RQ1)}
\label{sec:res-rq1}
% RESULTS
% P1 Answer: M0 higher on average in all protocols; all undetermined.
%    Gauged +0.020 [+0.007,+0.035], Holm p 0.053; random +0.015
%    [-0.010,+0.037], median +0.035, 66% of gauges, Holm p 0.012;
%    blocked +0.066 [+0.026,+0.105], 69%, Holm p 0.089 (Figure 5 row 1).
% P2 Lead: random +0.031 (24 h) to -0.003 (120 h); blocked +0.057 to +0.076.
%    40 leads: mean>0 at 36/13/40, band above 0 at 19/5/34 leads
%    (Figure 4 bottom).
% P3 Space (Figure 10): M0 higher at 70,75,79,65 of 98 gauges (6,12,18 h,
%    days 1-5, random) [add counts to facts sheet before citing].
%    By area (Figure 8c): equal < 50 km2; M0 higher > 600 km2.
% DISCUSSION
% P4 System difference (Section~\ref{sec:baselines}): upper bound on the
%    network's own contribution.
% P5 Literature: \citeA{Kirschstein2024} vs \citeA{WangH2025},
%    \citeA{Mosaffa2026}; our GEFS-driven, withheld-gauge setting.
% P6 Blocked: transfer to unseen terrain \cite{Feng2021}; 29 gauges in
%    3 basins and the LSTM's higher blocked KGE: suggestive only.

\begin{figure}
\figfile{figures/fig5_contrasts.pdf}{0.6\textwidth}
\caption{Paired differences in days 1--5 NSE for the four contrasts, by protocol: mean over gauges with 90\% gauge-bootstrap interval. The gray band marks the $\pm 0.02$ equivalence margin. M1 $-$ M4 is not available in the gauged protocol.}
\label{fig:contrasts}
\end{figure}

\begin{figure}
\figfile{figures/fig10_gauge_protocol_leads.pdf}{\textwidth}
\caption{Per-gauge difference in NSE, M0 $-$ LSTM, in the three protocols (rows) at leads of 6, 12, and 18~h and for days 1--5 (columns). Purple: M0 higher; orange: LSTM higher. The scale is clipped at $\pm 0.3$. Numbers give the gauges with M0 higher out of the gauges scored.}
\label{fig:maps-protocol}
\end{figure}

%---------------------------------------------------------------------
\subsection{Message Passing Versus Routing (RQ2)}
\label{sec:res-rq2}
% RESULTS
% P1 Message passing M0-M2: gauged +0.005 [+0.002,+0.009] equivalent
%    within 0.01; random +0.023 [-0.002,+0.054], median +0.001,
%    undetermined; blocked -0.013 [-0.038,+0.010] undetermined; sign
%    +,+,- across protocols.
% P2 Routing M1-M4: random +0.033 [+0.016,+0.053], median +0.026, 76% of
%    gauges, Holm p 1.4e-5: different. Positive at every lead 24-120 h
%    (+0.042 to +0.025); M1 higher at 73-81 of 98 gauges at 6-18 h
%    (Figure 9). By area: 600-2,000 km2 +0.055 [+0.026,+0.094] (86%),
%    < 200 km2 +0.005. Blocked -0.006 [-0.037,+0.021], undetermined; sign
%    changes at every lead.
% P3 Travel times: Spearman 0.88 with river length; median 10.3 h for
%    34.3 km; celerity 1.06 m/s. Peak timing (random) M1 3 h vs M4 6 h.
% DISCUSSION
% P4 Travel time, not neighbour mixing, is what gauges detect; consistent
%    with the gain peaking at intermediate basin sizes.
% P5 Message passing redundant with routed flow and broadcast specific
%    discharge; three-link reach; \cite{Alon2021,WangH2025}.
% P6 Differentiable routing \cite{Bindas2024,Song2025}; blocked transfer
%    of attribute-predicted routing parameters; n = 29.

\begin{figure}
\figfile{figures/fig9_gauge_contrast_leads.pdf}{\textwidth}
\caption{Per-gauge difference in NSE in the random hold-out ($n = 98$) for M0 $-$ LSTM, M0 $-$ M2 (message passing), M1 $-$ M4 (routing), and M0 $-$ M1 (insertion) (rows) at leads of 6, 12, and 18~h and for days 1--5 (columns). Purple: first model higher; orange: second model higher. The scale is clipped at $\pm 0.3$. Numbers give the gauges where the first model is higher out of the gauges scored.}
\label{fig:maps-contrasts}
\end{figure}

%---------------------------------------------------------------------
\subsection{The Value of Upstream Observations (RQ3)}
\label{sec:res-rq3}
% RESULTS
% P1 Answer: equivalent at days 1-5. Gauged +0.002 [+0.001,+0.004]
%    (equivalent within 0.01); random +0.000 [-0.007,+0.007] (equivalent
%    within 0.02). Blocked +0.012 is not a test (no visible upstream gauge):
%    noise floor, B1 +0.004, B2 +0.019, B3 +0.018.
% P2 Where insertion can act (random, 40 nested): 6 h +0.054
%    [+0.037,+0.072]; 12 h +0.009; 24 h -0.011; days 1-5 -0.003
%    [-0.018,+0.015]. Headwater about 0. Gauged nested 6 h +0.001
%    (Figure 7; s5_5a_position_by_lead.csv). Area-matched 12 pairs:
%    -0.0001 [-0.014,+0.013].
% P3 Exploratory (one sentence): area-share correlation -0.43 random,
%    +0.10 gauged; within-gauge test not identifiable (Text S/Figure S4).
% DISCUSSION
% P4 Other routes of observations (Section~\ref{sec:ablations}); M0-M1 is
%    the added value of the routed correction.
% P5 Median upstream travel time 10.3 h < first daily lead: consistent with
%    a 6 h-only gain. No-latency assumption makes it an upper bound.
% P6 Contrast with the expectation of Section~\ref{sec:obs}; practice:
%    insertion matters for nowcasting, not days 1-5.

\begin{figure}
\figfile{figures/fig7_upstream.pdf}{\textwidth}
\caption{Paired difference in NSE, M0 $-$ M1, against lead time at nested gauges with a visible upstream gauge and at headwater gauges in the (a) gauged and (b) random hold-out protocols: mean over gauges with 90\% gauge-bootstrap band.}
\label{fig:upstream}
\end{figure}

%---------------------------------------------------------------------
\subsection{Forecasts at Every Reach}
\label{sec:res-reach}
% RESULTS
% P1 Caveat first: reference = NWM v3.0 retrospective, M0's pretraining
%    target, never seen by the LSTM; agreement, not accuracy.
% P2 Median NSE by area (<5, 5-50, 50-500, 500-5,000, >5,000 km2):
%    M0 0.011, 0.112, 0.211, 0.306, -0.290; LSTM -0.723, 0.050, 0.197,
%    0.249, -1.449; NWM forecast -1.730, -0.696, -0.150, 0.211, -0.238.
%    Share NSE>0 (M0) 0.55, 0.81, 0.81, 0.72, 0.38. 529 of 1,569 reaches
%    > 5,000 km2 scored (Figure 8a,b; Figures S5).
% P3 Headwater reaches (23,172; median 1.96 km2): median NSE M0 0.014,
%    LSTM -0.650, NWM -1.643 (Figures S6).
% P4 Gauges by area (Figure 8c): 0.12 (< 50 km2, n = 3) to 0.48
%    (600-2,000 km2), random hold-out.
% DISCUSSION
% P5 Circularity: M0 inherits the reference at small streams.
% P6 Evidence gap below 50 km2 (77.7% of reaches; 4 gauges)
%    \cite{Krabbenhoft2022}.

\begin{figure}
\figfile{figures/fig8_skill_area.pdf}{\textwidth}
\caption{Skill by drainage area for M0, the LSTM, and the NWM forecast. (a) Median NSE, days 1--5, at every reach against the NWM v3.0 retrospective, on which M0 was pretrained; (b) share of those reaches with NSE $> 0$; (c) median NSE against observations at the evaluation gauges of the random hold-out. Numbers under the axis give the scored reaches (thousands) or gauges per class.}
\label{fig:area-skill}
\end{figure}

%---------------------------------------------------------------------
\subsection{Limitations}
\label{sec:limits}
% P1 One run per model and protocol; noise floor +0.012; +-0.02 margin
%    pre-chosen.
% P2 LSTM not input-matched; early checkpoints (13,754 of 68,767 updates).
% P3 NWM: nudged at withheld gauges; retrospective as pretraining target
%    and possible calibration at withheld gauges [VERIFY].
% P4 Blocked: three basins in the eastern highlands, n = 29.
% P5 Inputs: deterministic GEFS control member; no other forecast
%    variables; no data latency.
% P6 Coverage: no gauged M4; every-reach runs for M0 and LSTM only;
%    retrospective missing at most large reaches.

%---------------------------------------------------------------------
\subsection{What the River Network Contributes}
\label{sec:synthesis}
% One paragraph per RQ answer + one on implications for operational
% forecasting at unmonitored reaches (structure_md section 2, 5.7).

%=====================================================================
\section{Conclusions}
\label{sec:conclusions}
%=====================================================================
% One or two sentences per RQ with the headline number; what would change
% the answers (seeds, more blocked basins, small-stream gauges).
